\section{Percepción desde Silicon Valley: la ventana de progreso de los modelos se reordena}

La percepción de Guangmi (广密) es que, en el último trimestre, el nivel de los modelos avanzó más que en todo 2025; saltos como el de Opus 4.5 a 4.6 llevaron a los modelos de las preguntas y respuestas tipo chat a un modo verdaderamente agentic. La entrevista sitúa a las «tres grandes» de Silicon Valley en una ventana dinámica: OpenAI, Anthropic y Gemini tienen cada una su propia ventana de oportunidad, y la fórmula del éxito de hoy puede convertirse en el veneno de la siguiente etapa.

\begin{figure}[H]
\centering
\includegraphics[width=0.96\textwidth]{frontier-lab-map.png}
\caption{Mapa de posiciones de las empresas de modelos de Silicon Valley según las opiniones de la entrevista.}
\end{figure}

\begin{knowledgebox}{Lectura de la figura: no es una clasificación objetiva, sino un mapa de dependencia de la trayectoria}
La figura organiza las opiniones de la entrevista: Anthropic tiene una buena ventana en Coding/Agent; el éxito de OpenAI en ToC puede frenar su estrategia de Coding; Gemini tiene recursos y un ecosistema sólidos, pero le falta impulso; Meta se considera un retador; xAI se ve afectada por los vaivenes estratégicos y por problemas de equipo. El lector debe fijarse en la dependencia de la trayectoria, no tomarla como una clasificación fija.
\end{knowledgebox}

\subsection{Anthropic: detalles técnicos y ADN cultural}

En la entrevista se considera que Anthropic aprovechó la ventana de Coding/Agent. No tenía todo claro desde el day 1, pero sus fundadores revisan los datos de forma hands-on, dan importancia a los detalles técnicos y aplican entrevistas culturales estrictas; todo ello le facilita organizar recursos cuando surge un nuevo paradigma. Lo importante aquí es que la cultura no es un eslogan: influye en si el equipo valora cierto tipo de datos, de productos y de capacidades.

\subsection{OpenAI: las victorias pasadas pueden convertirse en veneno}

El éxito de OpenAI en ToC con ChatGPT podría, paradójicamente, llevar a la organización a concentrarse durante mucho tiempo en la puerta de entrada conversacional y en los productos de consumo, y a descuidar Coding. El juicio expresado en la entrevista es que la fórmula del éxito de una época puede convertirse en el veneno de la siguiente. Para una empresa de modelos, la inercia estratégica y los mecanismos de recompensa de la organización pueden ser más peligrosos que una carencia técnica puntual.

\subsection{Gemini, Meta y xAI}

Se valora que Gemini tiene modelos de gran capacidad y un buen nicho en el ecosistema, pero va muy rezagado en Coding; el error estratégico de Google fue no haber hecho de Coding su línea principal antes. Meta se considera un nuevo retador que sustituye a xAI como cuarto cabeza de serie; xAI, por su parte, se considera rezagada a corto plazo por sus vaivenes estratégicos y la salida de miembros de su founding team. Esto refleja que la competencia por los modelos de frontera ya no depende solo de los parámetros del modelo, sino de la organización, la estrategia, el producto y la capacidad de ejecución.

\begin{importantbox}{El verdadero problema de las tres grandes de Silicon Valley}
Una empresa de modelos líder que no dé importancia a Coding puede quedar fuera del primer nivel, porque Coding es la interfaz de ejecución con la que el modelo entra en el mundo digital y también la herramienta central para acelerar su propia investigación.
\end{importantbox}

\subsection{Resumen de la sección}

Las ventanas de las empresas de modelos rotan con rapidez. Quien logre convertir Coding/Agent en la línea principal de su organización podrá aprovechar el segundo acto; quien quede atrapado por el éxito de la etapa anterior podrá quedarse atrás.
